\documentclass[11pt]{article}

% ============================================================
%  RS Flavour Anarchy, the f-factor mechanism, and RS-GIM
%  The framework note that makes every flavour constraint
%  in our catalogue intelligible.
%  A slow, pedagogical internal review.
%  Written to track (a) the physics, (b) what is in OUR code,
%  (c) the subtleties and gaps.
%  This is an internal working note, not for distribution.
% ============================================================

\usepackage[margin=1in]{geometry}
\usepackage{amsmath,amssymb,amsthm}
\usepackage{mathtools}
\usepackage{graphicx}
\usepackage{booktabs}
\usepackage{enumitem}
\usepackage{xcolor}
\usepackage{framed}
\usepackage[colorlinks=true,linkcolor=blue,citecolor=blue,urlcolor=blue]{hyperref}

% lightweight "callout" box
\colorlet{shadecolor}{gray!12}
\newenvironment{callout}[1]{%
  \par\medskip\begin{shaded}\noindent\textbf{#1}\par\smallskip}{\end{shaded}\medskip}

\newcommand{\MKK}{M_{\mathrm{KK}}}
\newcommand{\LIR}{\Lambda_{\mathrm{IR}}}
\newcommand{\half}{\tfrac{1}{2}}
\newcommand{\vev}{v}
\newcommand{\eps}{\varepsilon}
\newcommand{\fIR}{f_{\mathrm{IR}}}
\newcommand{\fUV}{f_{\mathrm{UV}}}
\newcommand{\epsK}{\varepsilon_K}
\newcommand{\YEb}{\bar{Y}_E}
\newcommand{\YNb}{\bar{Y}_N}
\newcommand{\Mpl}{M_{\mathrm{Pl}}}
\newcommand{\gss}{g_{s*}}
\newcommand{\dmk}{\Delta m_K}
\newcommand{\code}[1]{\texttt{#1}}
\newcommand{\todo}[1]{{\color{red}[#1]}}

\theoremstyle{definition}
\newtheorem{definition}{Definition}

\title{\bf RS Flavour Anarchy, the $f$-Factor Mechanism, and RS-GIM\\[4pt]
\large A slow, pedagogical internal review --- physics, code, and subtleties}
\author{Internal working note (5D--Neutrino--Mixing repo)}
\date{\today}

\begin{document}
\maketitle

\begin{abstract}
\noindent
This is the \emph{framework} note: the conceptual root from which every flavour
constraint in our catalogue grows. It is a from-scratch, deliberately slow
walk-through of how a warped (Randall--Sundrum, RS) extra dimension turns
\emph{geometry} into the Standard-Model (SM) flavour puzzle's answer, and why that
same geometry, while spectacularly economical, is still not quite enough. It
assumes no prior knowledge of bulk fields, extra dimensions, or the flavour
problem: all are built up from the beginning. Part~\ref{part:loc} explains
fermion localisation --- the dimensionless bulk-mass parameter $c$, the
zero-mode profile, and the brane overlap factors $\fIR(c)$ and $\fUV(c)$ that our
code computes. Part~\ref{part:anarchy} shows how \emph{anarchic} (structureless,
$\mathcal{O}(1)$) five-dimensional Yukawa couplings, dressed by these overlap
factors, reproduce the hierarchical 4D charged-lepton, quark, and neutrino masses
\emph{and} the CKM/PMNS mixing angles --- with no hierarchical input couplings.
Part~\ref{part:rsgim} develops the RS-GIM mechanism: the automatic alignment that
suppresses flavour violation by the \emph{same} small overlaps that suppress the
masses, a built-in protection of GIM type. Part~\ref{part:gap} delivers the
punchline that makes the rest of the catalogue necessary --- why RS-GIM is
\emph{partial}, how the chirally-enhanced left--right $\Delta F=2$ operator slips
through it, and why $\epsK$ alone still demands $\MKK\sim10$--$20$~TeV or
alignment. Part~\ref{part:code} documents exactly what our code computes
(\code{warpConfig/wavefuncs.py}, the \code{yukawa} module) with every constant and
function named. Part~\ref{part:subtle} catalogues conventions, the universal-limit
caveat, and the honest scope. An annotated reading guide closes the note, and
cross-references point to the individual constraint notes ($\epsK$, $\Delta m_d$,
$\Delta m_s$, $D^0$, $\mu\to e\gamma$) that this framework underwrites.
\end{abstract}

\tableofcontents

% ====================================================================
\part*{Orientation: why this note exists}
\addcontentsline{toc}{section}{Orientation: why this note exists}
% ====================================================================

Our constraint catalogue is a collection of seemingly unrelated bounds: a
CP-violating kaon parameter $\epsK$, two $B$-meson mass splittings $\Delta m_d$
and $\Delta m_s$, charm mixing $D^0$--$\bar D^0$, the lepton-flavour-violating
decay $\mu\to e\gamma$, electroweak oblique parameters, $Z\to b\bar b$, and direct
collider searches. Each has its own note. But none of them can be \emph{read} ---
let alone trusted --- without first understanding the single mechanism they all
test. That mechanism is RS \emph{flavour anarchy}: the proposal that the
bewildering hierarchy of SM fermion masses and mixings is not built into the
fundamental couplings at all, but is generated by \emph{where each fermion lives}
in a warped fifth dimension. This note explains that idea from nothing, derives
the overlap factors our code actually evaluates, and then explains the two facts
that organise the entire catalogue: \emph{(a)} the geometry automatically aligns
flavour violation with the masses (RS-GIM), which is why the model is viable at
all, and \emph{(b)} the alignment is imperfect in one specific, chirally-enhanced
channel, which is why $\epsK$ and friends still bite. Everything downstream is a
quantitative working-out of those two sentences.

A one-paragraph executive summary, to be unpacked over the following pages:

\begin{quote}
\small
A 5D fermion in warped space is a field spread over the fifth dimension; a single
dimensionless number $c$ (its bulk mass in units of the curvature $k$) decides
whether its massless zero mode piles up near the Planck (UV) brane or near the
TeV (IR) brane. The SM Higgs lives on the IR brane, so a fermion's 4D mass is
proportional to how much of its zero mode reaches that brane --- the
\emph{overlap factor} $\fIR(c)$, which falls off exponentially as $c$ grows past
$\half$. Choosing the $c$'s to span a modest $\mathcal{O}(1)$ range turns
structureless, anarchic $\mathcal{O}(1)$ 5D Yukawas into the observed spread of
masses (the charged fermions alone, electron to top, already span more than five
orders of magnitude, and over twelve once the sub-eV neutrinos are included) and
the small CKM/PMNS mixing angles --- hierarchy from geometry, not from tuned
couplings. The same
overlaps that make a light fermion light also make its coupling to the
flavour-violating KK gauge bosons small, so flavour-changing amplitudes inherit
the mass suppression automatically: this is \emph{RS-GIM}, a built-in GIM-like
protection. It is enough to save the model in most channels --- but not all. The
chirality-flipping left--right four-quark operator in the kaon is enhanced by a
large hadronic chiral factor and by QCD running, and its surviving residue is
generically complex with an $\mathcal{O}(1)$ CP phase. RS-GIM suppresses the
coupling; the chiral and running enhancements give much of it back. The net is
the famous residual $\epsK$ problem: a generic anarchic model needs
$\MKK\sim10$--$20$~TeV, far above the TeV scale that motivated the construction,
unless one abandons pure anarchy for flavour \emph{alignment}.
\end{quote}

\noindent The logical spine of the whole note, in one picture:
\begin{footnotesize}
\begin{verbatim}
  bulk fermion in AdS5  ==>  one parameter c = M5/k  ==>  WHERE the zero mode lives
                          |
       c > 1/2:  UV-localised (Planck brane)   ==>  small IR overlap  ==>  LIGHT fermion
       c < 1/2:  IR-localised (TeV brane)      ==>  O(1) IR overlap   ==>  HEAVY fermion
                          |
  IR-brane Higgs samples the zero mode at the IR brane ==> overlap factor f_IR(c)
                          |
  m_fermion ~ 2 v k * f_L * Y5D * f_R     (Y5D anarchic, O(1), NO hierarchy)
  V_CKM,PMNS mixing ~ ratios of f-factors                  ==>  HIERARCHY FROM GEOMETRY
                          |
  KK gauge boson couples to fermion i with strength ~ g* f_i  (SAME f's as the mass)
                          |
       +-- RS-GIM:  flavour-violating coupling ~ g* f_i f_j  (light => small)
       |            ==> flavour violation ALIGNED with masses  ==  built-in GIM protection
       |
       +-- BUT partial:  surviving residue is COMPLEX, O(1) phase
                |
                +-- chirality-flipping LR operator: chiral factor ~26 x running ~8 = ~100
                |
  ==>  Im part of LR coefficient in the KAON  ==>  eps_K floor ~ 10-20 TeV
                          |
  custodial SU(2)_R x P_LR  ==>  fixes EW T and Z b_L b_L   ==  does NOTHING for eps_K
                          |
  what helps eps_K:   ALIGNMENT  |  larger bulk masses (c)  |  flavour symmetry
\end{verbatim}
\end{footnotesize}
\noindent Keep this map in mind; every box in it is unpacked below. The single
most common confusion --- that anarchy and RS-GIM together solve the RS flavour
problem outright --- is already half-defused here: the protection is real but
\emph{partial}, and the leak is the chirally-enhanced LR operator.

% ====================================================================
\part{Fermion localisation: the parameter $c$ and the overlap factors}
\label{part:loc}
% ====================================================================

\section{The warped geometry, in just enough detail}
\label{sec:geom}

Spacetime is a slice of five-dimensional anti-de~Sitter space (AdS$_5$) bounded by
two flat 4D branes. In the conformal coordinate $z$ used by the repo, the metric is
\begin{equation}
ds^2 = \left(\frac{1}{k z}\right)^2\!\big(\eta_{\mu\nu}\,dx^\mu dx^\nu - dz^2\big),
\qquad \frac{1}{k}\le z\le \frac{1}{\LIR},
\end{equation}
with $k$ the AdS curvature scale ($\sim\Mpl$). The \textbf{UV brane} sits at
$z_h=1/k$ (Planck-scale physics) and the \textbf{IR brane} at $z_v=1/\LIR$
(TeV-scale physics). The single small number that does all the work is the
\emph{warp factor}
\begin{equation}
\eps \;=\; \frac{\LIR}{k} \;=\; e^{-k\pi r_c}\;\sim\;10^{-15},
\label{eq:epsdef}
\end{equation}
where $r_c$ is the compactification radius. In code this is exactly
\code{epsilon = Lambda\_IR / k} (\code{warpConfig/baseParams.py}, function
\code{get\_warp\_params}), with the default $k=\Mpl=1.2209\times10^{19}$~GeV
(\code{MPL}) and a default $\LIR=3000$~GeV (\code{DEFAULT\_LAMBDA\_IR}), giving
$\eps\approx 2.5\times10^{-16}$. The exponentially small $\eps$ is what redshifts
Planck-scale physics down to the TeV scale; it is also, as we will see repeatedly,
the lever that converts $\mathcal{O}(1)$ inputs into enormous hierarchies.

\section{A bulk fermion and its zero mode}
\label{sec:zeromode}

A 5D Dirac fermion is not a point on a brane: it is a field $\Psi(x,z)$ spread
over the whole interval. Its 5D action contains a bulk Dirac mass which, by the
symmetries of AdS, is naturally measured in units of the curvature. We write that
dimensionless bulk mass as
\begin{equation}
\boxed{\;c \;\equiv\; \frac{M_5}{k}\;,\;}
\label{eq:cdef}
\end{equation}
the convention fixed in \code{CLAUDE.md} (so $M_5 = k\,c$, and the Bessel order
that controls the excited KK modes is $\alpha=|c+\half|$). Compactifying on the
interval, $\Psi$ decomposes into a Kaluza--Klein (KK) tower: one massless
\emph{zero mode} (which we identify with an SM fermion) plus an infinite tower of
massive excitations. Only the zero mode matters for the masses; the tower matters
for the constraints.

The crucial physics is \emph{where the zero mode lives}. Solving the zero-mode
equation in the warped background (Gherghetta--Pomarol, hep-ph/0003129;
Grossman--Neubert, hep-ph/9912408) gives a profile that is a power of the warp
coordinate, $\psi_0(z)\propto z^{2-c}$ up to normalisation. The exponent flips
character at $c=\half$:
\begin{itemize}[leftmargin=1.4em]
\item \textbf{$c>\half$}: the zero mode is peaked toward the \emph{UV} (Planck)
brane. Its tail reaching the IR brane is exponentially small. This is what we want
for a \emph{light} SM fermion.
\item \textbf{$c<\half$}: the zero mode is peaked toward the \emph{IR} (TeV)
brane, with an $\mathcal{O}(1)$ value there. This gives a \emph{heavy} fermion
(the top quark is the canonical IR-localised case).
\item \textbf{$c=\half$}: the conformal point --- the zero mode is flat, with no
preferred brane.
\end{itemize}
\textbf{Slogan to remember:} \emph{$c$ is a dial that slides a fermion between the
two branes; $c>\half$ pushes it to the Planck brane and makes it light, $c<\half$
pulls it to the TeV brane and makes it heavy.} Every mass and every mixing angle
in the model is, ultimately, a statement about a few $c$'s.

\section{The overlap factors $f_{\mathrm{IR}}(c)$ and $f_{\mathrm{UV}}(c)$}
\label{sec:ffactors}

The Higgs lives on (or very near) the IR brane. A fermion's 4D Yukawa coupling ---
and hence its mass --- is therefore proportional to \emph{how much of its
canonically-normalised zero mode reaches the IR brane}. That number is the
\textbf{IR overlap factor} $\fIR(c)$, the single most important function in the
whole framework. Normalising the profile $\psi_0\propto z^{2-c}$ over the warped
interval and evaluating at the IR brane gives, in the repo's convention,
\begin{equation}
\boxed{\;\fIR^2(c) \;=\; \frac{\half - c}{\,1 - \eps^{\,1-2c}\,}\;.}
\label{eq:fIR}
\end{equation}
For right-handed states that instead couple to UV-brane physics (the canonical
example is a UV-localised Majorana mass for a right-handed neutrino) the relevant
object is the \textbf{UV overlap factor}
\begin{equation}
\boxed{\;\fUV^2(c) \;=\; \frac{\half - c}{\,\eps^{\,2c-1} - 1\,}\;.}
\label{eq:fUV}
\end{equation}
These are exactly the two functions coded in \code{warpConfig/wavefuncs.py} as
\code{f\_IR(c, epsilon)} and \code{f\_UV(c, epsilon)} (both return $\sqrt{\cdot}$
of the above, vectorised over $c$, with the $c=\half$ limit handled separately ---
see \S\ref{sec:codeff}). They are the bridge between geometry and phenomenology.

It pays to read off the limiting behaviour, because the entire mass hierarchy is
hiding in it. For a \emph{UV-localised} fermion ($c>\half$, so $1-2c<0$), the term
$\eps^{1-2c}$ is a huge number ($\eps$ small to a negative power), and
\eqref{eq:fIR} collapses to
\begin{equation}
\fIR(c)\;\simeq\;\sqrt{c-\half}\;\;\eps^{\,c-\half}\qquad (c>\half),
\label{eq:fIRlimit}
\end{equation}
an \emph{exponentially small} number controlled by how far $c$ exceeds $\half$.
For an \emph{IR-localised} fermion ($c<\half$), the denominator of \eqref{eq:fIR}
is $\approx 1$ and $\fIR(c)\simeq\sqrt{\half-c}=\mathcal{O}(1)$. So:
\begin{center}
\begin{tabular}{lll}
\toprule
Regime & Overlap & SM interpretation \\
\midrule
$c>\half$ (UV-localised) & $\fIR\simeq\sqrt{c-\half}\,\eps^{c-\half}$ (exp.\ small) & light fermion \\
$c=\half$ (flat) & $\fIR^2 = 1/(-2\ln\eps) = 1/(2L)$ & intermediate \\
$c<\half$ (IR-localised) & $\fIR\simeq\sqrt{\half-c}=\mathcal{O}(1)$ & heavy fermion \\
\bottomrule
\end{tabular}
\end{center}
The middle row is the conformal point, where \eqref{eq:fIR} is singular ($0/0$)
and one takes the limit: $\fIR^2\to 1/(-2\ln\eps)$, which is $1/(2L)$ with
$L\equiv\ln(1/\eps)=k\pi r_c\approx 35$ the warp/volume log. This is precisely the
limit our code implements at $c=\half$ (\S\ref{sec:codeff}).

\begin{callout}{Worked numerical example --- reading the hierarchy off three $c$'s.}
Take the repo's validation geometry $\LIR=3000$~GeV, $k=\Mpl$, so
$\eps=\LIR/k\approx2.46\times10^{-16}$ and $L=\ln(1/\eps)\approx35.9$. Evaluate
$\fIR$ from \eqref{eq:fIR} at three illustrative bulk masses:
\begin{align*}
c=0.75:&\quad \fIR \approx \sqrt{0.25}\;\eps^{0.25}\approx 0.5\times(2.46\times10^{-16})^{0.25}\approx 6.3\times10^{-5}
   &&(\text{very UV, very light}),\\
c=0.60:&\quad \fIR \approx \sqrt{0.10}\;\eps^{0.10}\approx 0.32\times 0.027 \approx 8.7\times10^{-3}
   &&(\text{UV, light--ish}),\\
c=0.50:&\quad \fIR \approx \sqrt{1/(2L)}\approx \sqrt{1/71.9}\approx 0.12
   &&(\text{flat, intermediate}).
\end{align*}
These are the \code{c\_E = [0.75, 0.60, 0.50]} of the repo's benchmark
(\code{compute\_all\_yukawas} docstring), and they span more than \emph{three
orders of magnitude} in $\fIR$ (a factor $\sim1.9\times10^{3}$, from
$6.3\times10^{-5}$ to $0.12$) from an $\mathcal{O}(1)$ spread in $c$. That single fact ---
small change in $c$, exponential change in $f$ --- is the engine of the entire
construction. The IR overlap of the lepton doublet at $c_L=0.58$ is
$f_L\approx\sqrt{0.08}\,\eps^{0.08}\approx 0.016$, matching the docstring's
``\code{Should be \~0.016}''.
\end{callout}

% ====================================================================
\part{Anarchy: hierarchies and mixing without hierarchical couplings}
\label{part:anarchy}
% ====================================================================

\section{The flavour puzzle the mechanism solves}
\label{sec:puzzle}

Before showing the solution, state the problem. The SM charged fermion masses span
more than five orders of magnitude on their own --- $m_e\approx0.511$~MeV up to
$m_t\approx173$~GeV is a factor $\approx3.4\times10^{5}$ --- and the full fermion
spectrum reaches roughly twelve orders of magnitude once the sub-eV neutrinos are
counted
(in the lepton sector the repo stores $m_e=0.511$~MeV, $m_\mu=105.66$~MeV,
$m_\tau=1.777$~GeV in \code{yukawa/constants.py} as \code{M\_ELECTRON},
\code{M\_MUON}, \code{M\_TAU}). The CKM matrix is nearly the identity with small
off-diagonal angles; the PMNS matrix is the opposite, nearly anarchic with large
angles. In the SM these patterns are encoded in the Yukawa matrices by hand: one
simply \emph{chooses} couplings spanning this full range --- five-plus orders in
the charged sector, twelve once the neutrino Dirac couplings are included. That is
not an explanation; it is a restatement. The RS proposal is to explain the pattern
geometrically.

\section{Anarchic 5D Yukawas, dressed by overlaps}
\label{sec:anarchicmass}

The central hypothesis is \textbf{anarchy}: the fundamental 5D Yukawa couplings
$Y_5$ are \emph{structureless} --- complex numbers all of the same
$\mathcal{O}(1)$ order, with $\mathcal{O}(1)$ phases, and no hierarchy whatsoever.
All the structure comes from the overlap factors. For a charged lepton (Higgs and
Yukawa on the IR brane) the 4D mass is
\begin{equation}
\boxed{\;m_{E_i} \;=\; 2\,\vev\,k\;f_{L}\;Y_{E_i}\;f_{E_i}\;,}
\label{eq:massform}
\end{equation}
with $\vev\approx174$~GeV the electroweak vev (the repo's \code{V\_EWSB}), $k$ the
curvature, $f_L=\fIR(c_L)$ the doublet overlap and $f_{E_i}=\fIR(c_{E_i})$ the
right-handed overlap of generation $i$. This is exactly the formula in
\code{CLAUDE.md} and in \code{yukawa/charged\_lepton.py}. Read it as a product of
three pieces: a geometric prefactor $2vk$, a structureless anarchic Yukawa
$Y_{E_i}$, and the \emph{two} small overlaps that sandwich it. A light fermion is
light because \emph{both} its chiralities are UV-localised, so $f_L f_{E_i}$ is
tiny --- not because $Y_{E_i}$ is tiny.

The same statement for quarks (Gherghetta--Pomarol; Agashe--Perez--Soni
hep-ph/0408134) is the parametric
\begin{equation}
m_{u_i}\;\sim\;\vev\,(Y_u)_{ij}\,f_{Q_i}f_{u_j},\qquad
m_{d_i}\;\sim\;\vev\,(Y_d)_{ij}\,f_{Q_i}f_{d_j},
\label{eq:quarkmass}
\end{equation}
with anarchic $\mathcal{O}(1)$ matrices $(Y_{u,d})_{ij}$. The hierarchy of up-type
masses is the hierarchy of $f_{Q_i}f_{u_j}$; the top is heavy because both $f_{Q_3}$
and $f_{u_3}$ are $\mathcal{O}(1)$ (third generation IR-localised), the up quark
light because both its overlaps are exponentially small. \textbf{In one line:}
\emph{the many orders of magnitude in fermion masses --- five-plus across the
charged sector, twelve once neutrinos are folded in --- are exactly that many
orders of magnitude in $\prod f$, manufactured from an $\mathcal{O}(1)$ spread in
$c$, with the Yukawas doing none of the work.}

\section{Mixing angles for free}
\label{sec:mixing}

The most elegant part of the story is that the \emph{mixing} comes out right
automatically, with no extra input. Diagonalising \eqref{eq:quarkmass}, the
left-handed rotation that defines the CKM matrix has off-diagonal entries set by
ratios of the \emph{doublet} overlaps,
\begin{equation}
\boxed{\;|(V_{\rm CKM})_{ij}| \;\sim\; \frac{f_{Q_i}}{f_{Q_j}}\quad (i<j)\;,}
\label{eq:ckm}
\end{equation}
the parametric scaling stated by Agashe--Perez--Soni and Gherghetta--Pomarol. The
anarchic $\mathcal{O}(1)$ Yukawas drop out of the \emph{ratio}; only the geometry
survives. Plugging in the $c_Q$'s that already reproduce the up-type and down-type
masses, \eqref{eq:ckm} delivers the observed Cabibbo-sized hierarchy of CKM angles
--- a genuine \emph{postdiction}, not a fit. The same machinery in the lepton
sector gives the PMNS matrix; there the large measured angles are accommodated by
the more democratic (less hierarchical) lepton localisations, especially once the
seesaw is folded in.

\section{Neutrinos and the seesaw}
\label{sec:seesaw}

Neutrinos add one twist: a UV-localised right-handed neutrino $N$ with a Majorana
mass $M_N$ generated near the Planck brane. The light-neutrino mass then comes from
the type-I seesaw. In the universal limit (all $c_L$, $c_N$, $M_N$
generation-independent) the repo uses
\begin{equation}
m_{\nu_i} \;\simeq\; \frac{2\,k^2\,\vev^2\,f_L^2\,f_N^2}{(f_N^{\rm UV})^2\,M_N}\;Y_{N_i}^2,
\label{eq:seesaw}
\end{equation}
with $f_N=\fIR(c_N)$, $f_N^{\rm UV}=\fUV(c_N)$, and $Y_{N_i}$ the anarchic neutrino
Yukawa eigenvalues. This is the formula in \code{CLAUDE.md} and
\code{yukawa/neutrino.py}. Two features are worth flagging. First, the UV factor
$(f_N^{\rm UV})^{-2}$ is \emph{not} cosmetic: it is the canonical normalisation of
the UV-localised Majorana term, and using $\fUV$ \eqref{eq:fUV} rather than $\fIR$
here is the whole reason \code{wavefuncs.py} exports two functions. Second, the
seesaw $1/M_N$ explains the smallness of neutrino masses \emph{on top of} the
geometric suppression --- which is why neutrinos can be light without the
charged-lepton-style extreme UV localisation. The downstream $\mu\to e\gamma$ note
uses precisely the rescaled neutrino Yukawa as its flavour-violating spurion (see
\S\ref{sec:xref}).

\begin{callout}{Worked example --- inverting the formulas (the code's actual job).}
The code does not predict masses from $c$'s; it runs the logic \emph{backward} to
extract the anarchic Yukawas a given geometry would need, then checks they are
$\mathcal{O}(1)$ (perturbative). For charged leptons,
\code{compute\_charged\_lepton\_yukawas} inverts \eqref{eq:massform} to
\[
Y_{E_i} = \frac{m_{E_i}}{2\,\vev\,k\,f_L\,f_{E_i}},
\qquad
\bar{Y}_{E_i} \equiv 2k\,Y_{E_i} = \frac{m_{E_i}}{\vev\,f_L\,f_{E_i}}.
\]
The dimensionless \emph{rescaled} Yukawa $\bar Y=2kY$ is the physically meaningful
object: anarchy requires $|\bar Y|=\mathcal{O}(1)$, and the code's perturbativity
gate is $|\bar Y|<4$ (\code{is\_perturbative}, default \code{max\_Y\_bar = 4.0}).
At the benchmark ($\LIR=3000$, $c_L=0.58$, $c_E=[0.75,0.60,0.50]$) the docstring
notes $\bar Y_E$ comes out ``\code{O(1) to O(few)}'' --- exactly the anarchy
hypothesis self-consistently realised: the \emph{measured} lepton masses, divided
by the \emph{geometric} overlaps, leave behind structureless $\mathcal{O}(1)$
couplings. If instead they came out at $10^{-6}$ or at $10^3$, the geometry would
be doing too little or too much, and anarchy would be falsified for that point.
\end{callout}

% ====================================================================
\part{RS-GIM: flavour violation aligned with the masses}
\label{part:rsgim}
% ====================================================================

\section{Where the flavour violation comes from}
\label{sec:fcnc}

The same geometry that solves the mass puzzle creates a flavour problem, and the
two are inextricably linked --- which is the deep reason RS-GIM works. The SM gauge
bosons (gluon, $W$, $Z$, photon) are bulk fields; their flat zero modes are the
ordinary gauge bosons, but their \emph{KK excitations} are localised toward the IR
brane, with masses $m_n\simeq x_n\LIR$ set by Bessel roots (the first gauge KK root
for Neumann--Neumann boundary conditions is $x_1\approx2.45$; this $2.45$ is the
recurring $\MKK$-convention hazard of \S\ref{sec:mkkconv}). A KK gauge boson
couples to a fermion in proportion to their \emph{overlap} at the IR brane: a
heavy (IR-localised) fermion couples strongly, a light (UV-localised) fermion
weakly. Schematically the KK-gauge coupling to flavour $i$ is
\begin{equation}
g^{\rm KK}_{i}\;\sim\;\gss\,f_{q_i},
\label{eq:gKK}
\end{equation}
with $\gss$ the (volume-enhanced) 5D gauge coupling and $f_{q_i}$ the \emph{same}
overlap that set the mass. Because this coupling is \emph{not} flavour-universal
(it depends on $f_{q_i}$, which differs by generation), rotating to the fermion
mass basis does not diagonalise it: off-diagonal, flavour-changing KK-gauge
couplings survive. This is the origin of every $\Delta F=2$ and $\Delta F=1$
constraint in the catalogue.

\section{The RS-GIM mechanism}
\label{sec:rsgimmech}

Here is the saving grace. The off-diagonal (flavour-changing) KK-gauge coupling
between generations $i$ and $j$ is the product of two factors of \eqref{eq:gKK},
\begin{equation}
\boxed{\;(g^{\rm KK})_{ij}\;\sim\;\gss\,f_{q_i}\,f_{q_j}\;,}
\label{eq:rsgim}
\end{equation}
and those are the \emph{very same} overlap factors that suppress the masses of
generations $i$ and $j$. A flavour transition between two \emph{light} quarks
(say $s\to d$) is therefore automatically suppressed by $f_{q_1}f_{q_2}$, which is
small precisely because the down and strange quarks are light. The geometry that
explains \emph{why} they are light \emph{simultaneously} suppresses the
flavour-changing coupling between them. This is the \textbf{RS-GIM mechanism}, and
the analogy to the SM's own Glashow--Iliopoulos--Maiani mechanism is exact in
spirit: in the SM, flavour-changing neutral currents are suppressed by small
quark-mass-ratio factors and the GIM cancellation; in RS, they are suppressed by
small overlap factors that \emph{are} the mass ratios. As Agashe--Perez--Soni put
it, flavour violation between a zero mode and the KK sector is always accompanied
by a factor of the (small) overlap, so the dangerous couplings are aligned with the
masses by construction.

\textbf{Slogan to remember:} \emph{RS-GIM is the statement that you cannot make a
fermion light without also making its flavour violation small --- the two are the
same $f$-factor.} This is why the model is viable at all: a generic 5D theory with
$\mathcal{O}(1)$ flavour violation at the TeV scale would be excluded by
many orders of magnitude; RS-GIM buys back most of that automatically.

\section{Why it is automatic, not imposed}
\label{sec:automatic}

It is worth dwelling on \emph{why} this is not a tuning. In a generic
extra-dimensional or composite model one would have to \emph{arrange} for the
flavour-violating couplings to be small --- an extra assumption. In RS one gets it
for free, because the localisation that the masses \emph{force} is the same
localisation that the flavour-violating couplings \emph{inherit}. There is exactly
one set of $c$'s; it is fixed (up to anarchic $\mathcal{O}(1)$ rotations) by the
masses and mixings; and once fixed, the flavour violation is determined --- there
is no freedom left to make it dangerous. The price of this rigidity is the flip
side: there is also no freedom to make the \emph{residual} violation vanish, which
is the entire content of Part~\ref{part:gap}.

% ====================================================================
\part{Why the elegance is still not enough: the residual $\epsK$ problem}
\label{part:gap}
% ====================================================================

\section{RS-GIM is partial, and the leak is chiral}
\label{sec:partial}

RS-GIM suppresses the flavour-violating \emph{coupling} \eqref{eq:rsgim}, but
``suppressed'' is not ``zero,'' and three facts conspire to make the residue
matter --- nowhere more than in the CP-violating part of kaon mixing, the
constraint that defines the RS flavour problem.

\paragraph{(i) The residue is complex with an $\mathcal{O}(1)$ phase.} Anarchy
means the 5D Yukawas are complex with generic phases. RS-GIM rescales the
\emph{magnitude} of the surviving flavour violation by $f_i f_j$, but it does
nothing to the \emph{phase}: the leftover $(g^{\rm KK})_{sd}$ is generically
complex with an $\mathcal{O}(1)$ CP phase. There is no extra mechanism forcing it
real or aligning it with the SM. This is fatal for $\epsK$, which sees only the
\emph{imaginary} part of the mixing amplitude (full development in the
\code{epsilon\_k\_review} note).

\paragraph{(ii) The chirality-flipping LR operator is hadronically enhanced.} A
KK gluon couples to both chiralities, so it generates not only the SM-like
left--left operator $Q_1^{\rm VLL}$ but also the chirality-flipping left--right
operators $Q_4^{\rm LR}$, $Q_5^{\rm LR}$. In the kaon the matrix element of the LR
operator is enhanced over the LL one by the \emph{chiral factor}
\begin{equation}
r_\chi=\left(\frac{m_K}{m_s+m_d}\right)^2\approx 26,
\label{eq:chiral}
\end{equation}
because the pseudoscalar density is tied by the equations of motion to
$m_K^2/(m_s+m_d)$, a large number for the light kaon.

\paragraph{(iii) QCD running enhances the LR operator further.} Evolving the
Wilson coefficients from the multi-TeV matching scale down to the hadronic
$\sim2$~GeV mixes and \emph{enhances} the LR pair by a further factor $\sim8$. The
combined LR enhancement is of order $r_\chi\times8\sim10^2$ over the VLL estimate.

\noindent Put together: RS-GIM makes the LR \emph{coupling} small, but the chiral
and running enhancements give back two orders of magnitude, and the surviving
piece is complex. The product overshoots the tiny $\epsK$ budget unless the KK
scale is pushed up.

\section{Where the $10$--$20$~TeV floor comes from}
\label{sec:floor}

Because every Wilson coefficient scales as $1/\MKK^2$, the only knob that
suppresses the residual $\epsK$ contribution is to raise $\MKK$. Working through
the chain --- RS-GIM coupling \eqref{eq:rsgim}, times the LR chiral$\times$running
enhancement, times the $1/\dmk$ amplifier in the indirect-CP-violation master
formula --- the generic anarchic model with $\mathcal{O}(1)$ complex Yukawas
requires a first KK-gluon mass of
\begin{equation}
\boxed{\;\MKK \;\gtrsim\; \mathcal{O}(10\text{--}20)~\text{TeV}\quad(\text{generic anarchic})\;,}
\label{eq:floor}
\end{equation}
the headline result of Agashe--Perez--Soni (hep-ph/0408134, ``$\gtrsim
\mathcal{O}(10)$~TeV'') and, more sharply, Csaki--Falkowski--Weiler (0804.1954,
$\sim21$~TeV generic, $\sim33$~TeV in pseudo-Goldstone composite-Higgs variants).
This is roughly an order of magnitude above the few-TeV scale that motivated the
warped construction in the first place --- a residual ``little hierarchy.'' It is
\emph{the} RS flavour problem, and our $\epsK$ note computes it rigorously per
point.

\begin{callout}{The one number that organises the whole catalogue.}
$\epsK$ is the binding constraint because it stacks \emph{three} amplifiers on the
same complex residue: the chiral factor ($\sim26$), QCD running ($\sim8$), and the
$1/\dmk$ magnifier in $|\epsK|\simeq(\kappa_\eps/\sqrt2\,\dmk)\,{\rm Im}\,M_{12}$.
The other $\Delta F=2$ constraints share the mechanism but lack one or more
amplifier: $\Delta m_d$ and $\Delta m_s$ (the $B$ systems) have a much smaller
chiral factor and constrain the magnitude rather than a small imaginary part;
$D^0$ mixing is in the up sector and is theory-limited. So the same RS-GIM-leaking
LR operator gives a \emph{tiered} set of bounds --- $\epsK\sim20$~TeV the
sharpest, the $B$ and $D$ systems milder --- all traceable to this one framework.
The individual notes work out each tier.
\end{callout}

\section{What does and does not help}
\label{sec:helps}

The crucial model-dependence point, repeated in every downstream note: the
\emph{electroweak} custodial cure does \emph{nothing} for $\epsK$.
\begin{itemize}[leftmargin=1.4em]
\item \textbf{Custodial $SU(2)_R$ and $P_{LR}$} (the subject of the
\code{stu\_review} and \code{z\_to\_bb\_review} notes) protect the electroweak $T$
parameter and the $Z\bar b_L b_L$ vertex --- both $\Delta F=0/1$ \emph{electroweak}
effects. They say nothing about the colour sector or about the misalignment between
the KK-gluon basis and the quark mass basis. The $\Delta F=2$ floor
\eqref{eq:floor} survives the custodial construction essentially unchanged.
\item \textbf{What helps} is action on the \emph{flavour} structure, not the gauge
structure: (a) \emph{alignment} of the down-Yukawa and down-bulk-mass matrices (a
warped minimal-flavour-violation, e.g.\ the Rattazzi--Zaffaroni / shining
constructions and the alignment of 0907.0474) directly shrinks the dangerous
$(g^{\rm KK})_{sd}$; (b) \emph{larger bulk masses} (pushing light-generation $c$'s
deeper into the UV) suppress the overlaps at the cost of some naturalness; (c) an
imposed \emph{bulk flavour symmetry} ($U(3)$ or $U(2)$ families) that enforces
degeneracies.
\end{itemize}
All three act on $\mathrm{Im}(g_L)_{sd}(g_R)_{sd}$, the actual quantity $\epsK$
constrains. None of them is custodial symmetry.

\begin{callout}{Common confusions, defused.}
\begin{itemize}[leftmargin=1.2em,nosep]
\item[\textbf{A.}] \emph{``Anarchy + RS-GIM solves the RS flavour problem.''} No
--- it solves \emph{most} of it (the model is viable, which a generic 5D theory is
not), but the chirally-enhanced complex LR residue in $\epsK$ survives and forces
$\MKK\sim10$--$20$~TeV. RS-GIM is \emph{partial} protection.
\item[\textbf{B.}] \emph{``Custodial symmetry is the cure for everything in RS.''}
No --- custodial $SU(2)_R/P_{LR}$ cure the \emph{electroweak} $T$ and $Z\bar b_Lb_L$.
They do not touch the $\Delta F=2$ KK-gluon floor. Flavour needs alignment.
\item[\textbf{C.}] \emph{``$c$ is the fermion's 4D mass.''} No --- $c=M_5/k$ is a
\emph{5D bulk mass} in curvature units; it sets the zero-mode \emph{localisation},
and the 4D mass is $\propto f_L\,Y\,f_R$ with $f$ an exponential function of $c$.
A small change in $c$ is an exponential change in mass.
\item[\textbf{D.}] \emph{``Bigger $c$ means heavier.''} The opposite for the
\emph{zero mode}: $c>\half$ is UV-localised and \emph{light}; $c<\half$ is
IR-localised and \emph{heavy}. The dial runs backward from naive intuition.
\item[\textbf{E.}] \emph{``$\MKK$ is one number.''} It is overloaded by
$x_1\approx2.45$: the physical first KK mass is $x_1\LIR$, the repo bookkeeping
default is $\MKK=\LIR$. A ``$20$~TeV'' bound must say which (\S\ref{sec:mkkconv}).
\end{itemize}
\end{callout}

% ====================================================================
\part{What is in OUR code}
\label{part:code}
% ====================================================================

\emph{Reader here for the physics may skim Part~\ref{part:code} on a first pass
and jump to the subtleties in Part~\ref{part:subtle}; this part documents the
exact code so the note doubles as an implementation reference. Unlike the
constraint notes, this framework module contains no veto: it computes the
\emph{inputs} (overlaps and anarchic Yukawas) that the constraint modules
consume.}

\section{The overlap factors: \code{warpConfig/wavefuncs.py}}
\label{sec:codeff}

The two functions are \code{f\_IR(c, epsilon)} and \code{f\_UV(c, epsilon)}, both
returning the square root of \eqref{eq:fIR} and \eqref{eq:fUV} respectively, both
vectorised over array-valued $c$. The exact implemented expressions are
\begin{align}
\code{f\_IR}:\quad \fIR^2 &= \frac{0.5-c}{1-\eps^{\,1-2c}}, &
\code{f\_UV}:\quad \fUV^2 &= \frac{0.5-c}{\eps^{\,2c-1}-1}.
\end{align}
Three implementation details a reader of the code should know:
\begin{itemize}[leftmargin=1.4em]
\item \textbf{The $c=\half$ limit is handled separately.} A boolean mask
\code{mask = \~np.isclose(c\_arr, 0.5)} routes $c=\half$ to the analytic limit
\code{res\_sq[\~mask] = 1.0 / (-2.0 * np.log(epsilon))}, i.e.\ $\fIR^2(\half)=1/(2L)$,
avoiding the $0/0$ in the closed form. Both functions use the identical limit (it
is the same value for IR and UV at the conformal point).
\item \textbf{Negative-noise guard.} \code{res\_sq = np.maximum(res\_sq, 0.0)}
before the square root protects against tiny negative values from floating-point
cancellation. So the functions never return \code{nan} from a near-zero negative.
\item \textbf{No veto here.} \code{wavefuncs.py} has \emph{no} severity, anchor, or
budget --- it is a pure geometry primitive consumed by everything downstream.
\end{itemize}

\section{The geometry: \code{warpConfig/baseParams.py}}
\label{sec:codegeom}

\code{get\_warp\_params(k, Lambda\_IR)} returns the dictionary
$\{k,\LIR,\eps,r_c,z_h,z_v,\text{warp\_log}\}$ with
\code{epsilon = Lambda\_IR / k}, \code{warp\_log = -log(epsilon)} (this is
$L=\ln(1/\eps)$), and \code{rc = warp\_log / (pi * k)}. The module constants are
\code{V\_EWSB = 174.0}~GeV, \code{MPL = 1.2209e19}~GeV,
\code{MBARPL = 2.435e18}~GeV (reduced Planck mass),
\code{DEFAULT\_K = MPL}, and \code{DEFAULT\_LAMBDA\_IR = 3000.0}~GeV. The function
raises \code{ValueError} if $\eps\le0$. Note the repo's $\LIR$ is the
\emph{geometric} IR scale $1/z_v$, \emph{not} the physical first KK mass --- the
$x_1\approx2.45$ distinction of \S\ref{sec:mkkconv}.

\section{The Yukawa module: inverting the mass formulas}
\label{sec:codeyuk}

The \code{yukawa} package turns a geometry point into the anarchic Yukawas it
implies, and checks perturbativity. The entry point
\code{compute\_all\_yukawas(Lambda\_IR, c\_L, c\_E, c\_N, M\_N,
lightest\_nu\_mass, ordering, ...)} (\code{yukawa/compute\_yukawas.py}) does, in
order:
\begin{enumerate}[leftmargin=1.6em]
\item geometry via \code{get\_warp\_params}, then the overlaps
\code{f\_L = f\_IR(c\_L, eps)}, \code{f\_E = f\_IR(c\_E\_arr, eps)},
\code{f\_N = f\_IR(c\_N, eps)}, \code{f\_N\_UV = f\_UV(c\_N, eps)};
\item charged-lepton Yukawas via \code{compute\_charged\_lepton\_yukawas}
(\code{yukawa/charged\_lepton.py}), inverting \eqref{eq:massform}:
\code{Y\_E = m\_E / (2.0 * v * k * f\_L * f\_E\_arr)} and the rescaled
\code{Y\_E\_bar = 2.0 * k * Y\_E};
\item neutrino Yukawas via \code{compute\_neutrino\_yukawas}
(\code{yukawa/neutrino.py}), inverting the seesaw \eqref{eq:seesaw} with
\code{prefactor = (f\_N\_UV**2 * M\_N) / (2.0 * k**2 * v**2 * f\_L**2 * f\_N**2)}
and \code{Y\_N = sqrt(m\_nu\_GeV * prefactor)};
\item the PMNS-dressed neutrino matrix
\code{Y\_N\_matrix = V\_pmns @ diag(Y\_N)}, encoding $Y_N\to V_{\rm PMNS}\,{\rm
diag}(Y_{N_i})$ in the charged-lepton mass basis.
\end{enumerate}
The lepton masses are the PDG values \code{LEPTON\_MASSES = (M\_ELECTRON,
M\_MUON, M\_TAU)} from \code{yukawa/constants.py} (\code{M\_ELECTRON =
0.51099895e-3}, \code{M\_MUON = 105.6583755e-3}, \code{M\_TAU = 1.77686}~GeV);
neutrino masses are converted eV$\to$GeV with \code{EV\_TO\_GEV = 1e-9}.

\section{The perturbativity gate}
\label{sec:codepert}

The result object \code{YukawaResult} exposes
\code{is\_perturbative(max\_Y\_bar=4.0)}, returning \code{True} iff
\emph{all} $|\bar Y_E|$ and $|\bar Y_N|$ are below $4.0$. The docstring explains
the choice: $|\bar Y|<4$ corresponds to ``at least $\sim3$ KK modes before the
Yukawa becomes strongly coupled.'' This is the \emph{anarchy self-consistency
check}: if the rescaled Yukawas a geometry implies are $\mathcal{O}(1)$ (and
below~4), the point realises anarchy honestly; if they blow up, the geometry is
over-stretched. Input validation in both inverters raises \code{ValueError} on
non-positive $f$-factors or $M_N$, or wrong-shape inputs.

\section{How a downstream constraint consumes this}
\label{sec:xref}

This module is the \emph{root}, not a leaf, of the constraint graph. The
flavour-violating amplitudes the catalogue vetoes are built from the same
$f$-factors and (rescaled) Yukawas computed here:
\begin{center}
\begin{tabular}{lll}
\toprule
Downstream note / constraint & What it pulls from this framework & Severity \\
\midrule
\code{epsilon\_k\_review} / \code{K001} & RS-GIM coupling $\gss f_i f_j$, anarchic complex phase & HARD (rigorous) \\
\code{delta\_m\_d\_review} / \code{B001} & same KK-gluon $\Delta F=2$, $B_d$ system & HARD \\
\code{delta\_m\_s\_review} / \code{B003}/\code{B004} & same, $B_s$ system ($\Delta m_s$, $\phi_s$) & HARD \\
\code{d0\_mixing\_review} / \code{C001}/\code{C002} & up-sector $\Delta F=2$, same mechanism & HARD magnitude; CP proxy \\
\code{mu\_e\_gamma\_review} / \code{L001} & rescaled $\bar Y_N$ spurion $\bar Y_N\bar Y_N^\dagger$ & HARD (proxy/NDA) \\
\bottomrule
\end{tabular}
\end{center}
In particular the $\mu\to e\gamma$ note's flavour-violating spurion is built from
exactly the \code{Y\_N\_bar} this module returns (the $\bar Y_N=2k\,Y_N$
rescaling is the repo-wide convention), and the $\Delta F=2$ notes all inherit the
RS-GIM coupling structure \eqref{eq:rsgim} derived here. Read this note first; the
others are its consequences.

% ====================================================================
\part{Subtleties and the gap to ``rigorous''}
\label{part:subtle}
% ====================================================================

\section{The universal-limit caveat in the seesaw}
\label{sec:univlimit}

The seesaw formula \eqref{eq:seesaw} as coded assumes the \emph{universal limit}:
$c_L$, $c_N$, and $M_N$ are generation-independent, so the neutrino mass formula
reduces to a scalar relation per eigenvalue. The neutrino docstring
(\code{yukawa/neutrino.py}) states this explicitly: ``\emph{The formula assumes the
universal limit. For non-universal bulk masses, the full matrix structure would
need to be considered.}'' A fully general treatment would carry the
$3\times3$ matrix structure of $f_L$, $f_N$ and the Majorana mass. For the
charged-lepton sector the $c_{E_i}$ \emph{are} generation-dependent (the
inversion is per-eigenvalue and exact for diagonal $Y_E$), but the off-diagonal
anarchic entries are not reconstructed by the scalar inverter --- only the
diagonal eigenvalues are.

\section{Benchmark is paper-inspired, not a literal reproduction}
\label{sec:benchmark}

The repeatedly-used validation point ($c_L=0.58$, $c_N=0.27$,
$c_E=[0.75,0.60,0.50]$, $\LIR=3000$~GeV, $M_N=1.22\times10^{18}$~GeV) is flagged
in both \code{compute\_yukawas.py} and \code{neutrino.py} as a \emph{repo-local,
paper-inspired benchmark}: ``\emph{It should not be described as an exact
reproduction of the paper's Eq.~(10) Yukawa values}'' (referencing Perez--Randall,
arXiv:0805.4652). It reproduces the right \emph{orders of magnitude} ($f_L\approx
0.016$, $\bar Y_E\sim\mathcal{O}(1\text{--few})$, $\bar Y_N\sim\mathcal{O}(0.2$--$1)$),
which is the point of anarchy, but is not a digit-for-digit table match. Treat the
benchmark numbers as illustrative, not as a literature anchor.

\section{The $M_{\rm KK}$ convention: the factor of $2.45$}
\label{sec:mkkconv}

The single most important bookkeeping hazard, shared with every downstream note.
The geometric IR scale $\LIR=1/z_v$ (the repo's \code{Lambda\_IR}) is \emph{not}
the physical first KK mass. The physical first KK gauge mass is
$\MKK^{\rm phys}=x_1\LIR$ with $x_1\approx2.45$ (the Neumann--Neumann root
recorded elsewhere in the repo as \code{GAUGE\_KK\_ROOT\_NN =
2.448687135269161}). A bound quoted as ``$20$~TeV'' in $\LIR$ units is
``$\sim49$~TeV'' in physical first-KK units; the literature's $\sim21$~TeV is a
\emph{physical} KK-gluon mass. When a downstream constraint forms $1/\MKK^2$, the
convention fed in must match the convention the bound was calibrated in --- mixing
them mis-scales the amplitude by $2.45^2\approx6$. This framework module itself
exposes only $\LIR$ (\code{baseParams.py}); the $\MKK$ convention is fixed
downstream.

\section{Anarchy is an assumption, not a derivation}
\label{sec:anarchyassumption}

The $\mathcal{O}(1)$-Yukawa anarchy hypothesis is an \emph{input}, not a theorem.
The code does not \emph{prove} the surviving Yukawas are anarchic; it
\emph{checks} (via \code{is\_perturbative}) that a given geometry is consistent
with anarchy. A point that needs $|\bar Y|>4$ is rejected as over-stretched, but a
point that passes is only \emph{compatible} with anarchy --- the actual phases and
$\mathcal{O}(1)$ magnitudes are not predicted. This is exactly why the residual
$\epsK$ bound is a \emph{floor on a flavour structure} (anarchy), not a law of all
warped models: imposing alignment instead (\S\ref{sec:helps}) changes the
surviving couplings and relaxes the bound.

\section{What this framework module does \emph{not} do}
\label{sec:scope}

To be complete about scope:
\begin{itemize}[leftmargin=1.4em]
\item \textbf{No KK-tower spectral computation.} \code{wavefuncs.py} gives the
\emph{zero-mode brane overlap} in closed form; it does not solve for the KK-tower
profiles or the $\gss$ KK-gauge couplings (those live in the quark-coupling
machinery the $\Delta F=2$ notes describe, and are themselves a model layer).
\item \textbf{Diagonal-Yukawa inversion only.} The Yukawa module reconstructs
\emph{eigenvalues}, not the full anarchic off-diagonal matrices. The mixing
\eqref{eq:ckm} is parametric framework physics, applied in
the constraint modules, not reconstructed here.
\item \textbf{No CP phases here.} The complex $\mathcal{O}(1)$ phases that make
$\epsK$ dangerous (\S\ref{sec:partial}) are a property of the anarchic Yukawas
assumed downstream; this module returns real-magnitude $f$-factors and
($\sqrt{\cdot}$ of real) Yukawa eigenvalues.
\end{itemize}

\section{Checklist: what a maximally rigorous framework layer would add}
\begin{enumerate}[leftmargin=1.8em]
\item Generalise the seesaw inverter to the non-universal $3\times3$ case
(\S\ref{sec:univlimit}).
\item Reconstruct the full anarchic Yukawa matrices (off-diagonal entries and
phases), not just eigenvalues, so the CKM/PMNS mixing \eqref{eq:ckm} is computed
rather than assumed.
\item Replace the closed-form brane overlaps with explicit tower-summed warped
profile integrals where sub-leading corrections matter.
\item Fix and document a single $\MKK$ convention threaded through every
downstream $1/\MKK^2$ (\S\ref{sec:mkkconv}).
\item Make the anarchy assumption falsifiable point-by-point by carrying the
implied Yukawa \emph{distribution}, not just the perturbativity flag.
\end{enumerate}

% ====================================================================
\part*{Annotated reading guide}
\addcontentsline{toc}{section}{Annotated reading guide}
% ====================================================================

\begin{description}[leftmargin=0em,style=nextline]
\item[Gherghetta, Pomarol, hep-ph/0003129 (\emph{Nucl.\ Phys.\ B586 (2000)
141}).] The foundational paper on bulk fermions in RS. The zero-mode profile
$\psi_0\propto z^{2-c}$, the $c\gtrless\half$ localisation dichotomy of
\S\ref{sec:zeromode}, and the mass scaling $m\sim Y_5\,k\,f_Q f_u$ of
\eqref{eq:quarkmass} are theirs. Read \S\ref{part:loc} alongside it.

\item[Grossman, Neubert, hep-ph/9912408.] Neutrinos in the bulk: the localisation
mechanism for right-handed neutrinos and the warped seesaw structure underlying
\eqref{eq:seesaw}. The reference for \S\ref{sec:seesaw} and the $\fUV$ factor.

\item[Huber, hep-ph/0303183.] Flavour in RS from anarchic 5D Yukawas: the explicit
demonstration that an $\mathcal{O}(1)$ spread of bulk masses reproduces the quark
masses \emph{and} the CKM angles \eqref{eq:ckm}. The quantitative backbone of
\S\ref{part:anarchy}.

\item[Agashe, Perez, Soni, hep-ph/0408134.] The flavour structure of warped
models and the RS-GIM mechanism: that flavour violation between a zero mode and the
KK sector is always accompanied by the small overlap factor (\S\ref{sec:rsgimmech}),
and that despite this protection the chirally-enhanced $\epsK$ contribution
survives and forces $\MKK\gtrsim\mathcal{O}(10)$~TeV. The reference for
\S\ref{part:rsgim} and \S\ref{part:gap}.

\item[Csaki, Falkowski, Weiler, arXiv:0804.1954 (\emph{JHEP 0809:008}).] The
sharpest statement of the residual problem: the KK-gluon $\Delta F=2$ Hamiltonian,
the chirality-flipping $C_4/C_5$ operators, the much tighter bound on the
\emph{imaginary} part, and the $\sim21$/$\sim33$~TeV floors. This is the paper our
$\Delta F=2$ core matches to machine precision (see the \code{epsilon\_k\_review}
note). Read \S\ref{part:gap} alongside it.

\item[Our companion notes (same repo, \code{review\_local/}).] The constraint
notes that this framework underwrites: \code{epsilon\_k\_review} (the sharpest
$\Delta F=2$ bound; the LR/chiral/running chain in full),
\code{delta\_m\_d\_review} and \code{delta\_m\_s\_review} (the $B$-system
companions), \code{d0\_mixing\_review} (the up-sector case), and
\code{mu\_e\_gamma\_review} (the lepton-sector RS-GIM and the $\bar Y_N$ spurion).
The electroweak-side notes \code{stu\_review} and \code{z\_to\_bb\_review} cover
the \emph{custodial} cures that, as \S\ref{sec:helps} stresses, do \emph{not}
touch the flavour floor. Start here, then branch.
\end{description}

\bigskip
\noindent\rule{\linewidth}{0.4pt}\\
\footnotesize
\emph{Internal note. The overlap formulas \eqref{eq:fIR}--\eqref{eq:fUV}, the mass
and seesaw relations \eqref{eq:massform}, \eqref{eq:quarkmass}, \eqref{eq:seesaw},
and the code expressions of Part~\ref{part:code} are exactly what
\code{warpConfig/wavefuncs.py}, \code{warpConfig/baseParams.py}, and the
\code{yukawa} module compute; the constants ($v=174$~GeV, $k=\Mpl=1.2209\times
10^{19}$~GeV, the PDG lepton masses, \code{max\_Y\_bar = 4.0}, the $c=\half$ limit
$1/(2L)$) are as stored in the repo. The RS-GIM, anarchy, and residual-$\epsK$
physics follow the cited papers; the ``rigorous vs proxy'' tags and the $\MKK$
convention follow the orchestration ledgers and the companion constraint notes.
Cross-check before citing externally.}

\end{document}
